\section{Full Proof of Proposition~\ref{prop:orthogonality}}
\label{sec:proof_orthogonality}

We provide a self-contained derivation of the bound $|\cosff| \lesssim 0$
under the two conditions stated in Proposition~\ref{prop:orthogonality}.
All quantities are in the $(A,B)$ parameter space $\mathbb{R}^p$ unless noted.

\subsection*{Step 1: Jacobian decomposition}

For each layer $l$ with $B_l\in\mathbb{R}^{d\times r}$, $A_l\in\mathbb{R}^{r\times d}$,
the Jacobian of $\mathrm{vec}(\DeltaW_l)$ with respect to $(A_l, B_l)$ is (Eq.~\eqref{eq:jacobian}):
\begin{equation}
  J_l = \bigl[I_d\otimes B_l \;\Big|\; A_l^\top\otimes I_d\bigr]
  \;\in\; \mathbb{R}^{d^2\times 2rd}.
\end{equation}
The full Jacobian $J\in\mathbb{R}^{Ld^2\times p}$ is block-diagonal in $l$.
Its Gram matrix factorizes as:
\begin{equation}
  J_l J_l^\top = (I_d\otimes B_lB_l^\top) + (A_lA_l^\top\otimes I_d).
  \label{eq:gram}
\end{equation}

\subsection*{Step 2: Gradient vectors in $(A,B)$-space}

\textbf{Flat gradient.}
Using the Gauss-Newton approximation $H_\phi \approx J^\top H_{\DeltaW} J$
(valid when task loss is small or at local optima):
\begin{equation}
  \ggam - \gclean \;\approx\; H_\phi\,\epsilon \;\approx\; J^\top H_{\DeltaW} J\,\epsilon,
  \label{eq:gflat_approx}
\end{equation}
where $\epsilon = \rho\,\gclean/\|\gclean\|$ is the perturbation direction.

\textbf{Fisher gradient.}
From Eq.~\eqref{eq:gfisher_chain}:
$\gfisher = J^\top \mathbf{f}$,
where $\mathbf{f} = \lambdaewc\,\mathrm{vec}(\Ftilde_l\odot d_l)$ with $d_l = \DeltaW_l - \DeltaW_l^*$.

\subsection*{Step 3: Inner product expansion}

\begin{align}
  (\ggam-\gclean)^\top\gfisher
  &\;\approx\; (J^\top H_{\DeltaW} J\epsilon)^\top (J^\top\mathbf{f}) \notag\\
  &\;=\; \epsilon^\top J^\top H_{\DeltaW}\,(JJ^\top)\,\mathbf{f}.
  \label{eq:inner_product}
\end{align}

\subsection*{Step 4: Applying condition (ii) — Jacobian isotropy}

Under condition~(ii), $JJ^\top \approx c^2\,I_{Ld^2}$, so:
\begin{equation}
  (\ggam-\gclean)^\top\gfisher
  \;\approx\; c^2\,\epsilon^\top J^\top (H_{\DeltaW}\,\mathbf{f}).
  \label{eq:after_isotropy}
\end{equation}
The Jacobian isotropy condition is satisfied when
$A_lA_l^\top\approx \sigma_A^2 I_d$ and $B_lB_l^\top\approx \sigma_B^2 I_d$
for each layer (Eq.~\eqref{eq:gram}).
This holds approximately when the LoRA factors have balanced singular values,
which is the generic case for matrices initialized by Gaussian entries and updated by SGD
with moderate step sizes.
Condition~(ii) fails only in degenerate regimes
(e.g., $r=1$ where $A_lA_l^\top$ has rank 1,
or if one factor collapses to near-zero).

\subsection*{Step 5: Applying condition (i) — Spectral separation}

Eq.~\eqref{eq:after_isotropy} contains the term $H_{\DeltaW}\,\mathbf{f}$,
where $\mathbf{f}$ lies in the column space of $F_{t'}$ (since $\mathbf{f} = \lambdaewc F_{t'}\odot d$,
and $F_{t'}$ acts as a diagonal scaling).

Let $U_H \in \mathbb{R}^{Ld^2\times k}$ be the top-$k$ eigenvectors of $H_{\DeltaW}$
and $U_F \in \mathbb{R}^{Ld^2\times k}$ the top-$k$ eigenvectors of $F_{t'}$.
The vector $\epsilon^\top J^\top(H_{\DeltaW}\,\mathbf{f})$ can be bounded by:
\begin{align}
  \bigl|\epsilon^\top J^\top H_{\DeltaW}\,\mathbf{f}\bigr|
  &= \bigl|(J\epsilon)^\top H_{\DeltaW}\,\mathbf{f}\bigr| \notag\\
  &\leq \|J\epsilon\|\,\|H_{\DeltaW}\|\,\|(I - U_H U_H^\top)\mathbf{f}\|
    + \|J\epsilon\|\,\|H_{\DeltaW}\|\,\|U_H^\top\mathbf{f}\|.
  \label{eq:bound}
\end{align}
Under condition~(i), the principal angles $\theta_i(U_H, U_F)$ between the two eigenspaces
are all large (near $\pi/2$), so the cross-overlap matrix satisfies:
\begin{equation}
  \|U_H^\top U_F\|_F \;=\; \Bigl(\sum_i \cos^2\theta_i(U_H,U_F)\Bigr)^{1/2} \;\approx\; 0.
  \label{eq:uh_uf_bound}
\end{equation}
(Note: $\|U_H^\top U_F\|_F$ depends on $\cos\theta_i$, not $\sin\theta_i$; when eigenspaces are
near-orthogonal, all $\cos\theta_i \approx 0$, driving the overlap to zero.)
Since $\mathbf{f} \approx U_F U_F^\top \mathbf{f}$:
\begin{equation*}
  \|U_H^\top\mathbf{f}\| \;\leq\; \|U_H^\top U_F\|_F\,\|U_F^\top\mathbf{f}\|
  \;\leq\; \|U_H^\top U_F\|_F\,\|\mathbf{f}\| \;\approx\; 0.
\end{equation*}

\subsection*{Step 6: Final bound}

Combining Steps 3--5:
\begin{equation}
  |\cosff| \;=\; \frac{|(\ggam-\gclean)^\top\gfisher|}{\|\ggam-\gclean\|\,\|\gfisher\|}
  \;\lesssim\;
  c^2\,\|H_{\DeltaW}\|\,\frac{\|J\epsilon\|}{\|\ggam-\gclean\|}
  \cdot\|U_H^\top U_F\|_F.
  \label{eq:final_bound}
\end{equation}
The factor $\|J\epsilon\|/\|\ggam-\gclean\|$ is $\mathcal{O}(1)$ (both are $\Theta(1)$ in the gradient norm).
Under condition~(i), $\|U_H^\top U_F\|_F = (\sum_i\cos^2\theta_i)^{1/2} \approx 0$
when all principal angles $\theta_i \approx \pi/2$.
Thus $|\cosff| = \mathcal{O}(\|U_H^\top U_F\|_F)$:
near-zero cosine follows directly from large principal angles between the current-task Hessian
and the past-task Fisher in $\DeltaW$-space, \emph{provided} that the Jacobian is approximately isotropic.

\subsection*{When condition (ii) is tight}

If condition~(ii) fails (e.g., $B_lB_l^\top$ is rank-$r$ concentrated on one direction),
then $JJ^\top$ has large eigenvalue spread,
and the intermediate step Eq.~\eqref{eq:after_isotropy} does not simplify to
a scalar times the identity.
In that case, the inner product in Eq.~\eqref{eq:inner_product} involves
$(JJ^\top)\mathbf{f}$, which may rotate $\mathbf{f}$ into the $H_{\DeltaW}$ eigenspace
even when $H_{\DeltaW}$ and $F_{t'}$ are spectrally separated.
The empirical observation that $\cosff\approx 0$ across all our experiments with $r\in\{10,16\}$
suggests that condition~(ii) holds throughout training for these rank choices.

\section{Per-Task Accuracy Matrices}
\label{sec:per_task}

\begin{table}[h]
\caption{Final task-level accuracies for FS-LoRA (ewc\_lambda=2000, seed 0).
Each row shows the accuracy on all tasks at the end of training the last task.}
\label{tab:per_task}
\centering
\begin{tabular}{lp{10cm}}
\toprule
Dataset & Final task accuracies \\
\midrule
CUB200    & [92.21, 78.15, 90.64, 79.67, 91.40, 76.13, 81.40, 83.04, 82.28, 92.24] \\
Cars196   & [80.50, 45.68, 61.06, 65.15, 60.67, 66.30, 64.22, 65.85, 66.13, 52.74] \\
Aircraft  & [53.15, 39.94, 40.12, 59.16, 71.77, 66.47, 69.37, 73.57, 60.18, 72.97] \\
Flowers   & [98.68, 98.04, 93.70, 86.58, 79.68, 93.46, 90.99, 90.95, 89.38, 82.97] \\
OxfordPet & [92.57, 90.18, 87.66, 97.25, 97.74, 95.00, 87.50, 97.24, 96.40] \\
\bottomrule
\end{tabular}
\end{table}

\section{Hyperparameter Sensitivity}
\label{sec:hparam}

\begin{table}[h]
\caption{$\lambdaewc$ sweep on CUB200 (trace-normalized, seed 0).
Results are stable across a wide range of $\lambdaewc$.}
\label{tab:lam_sweep}
\centering
\begin{tabular}{ccccc}
\toprule
$\lambdaewc$ & CNN FAA & CNN AAA & Forget & BWT \\
\midrule
500  & 73.931 & 74.241 & 11.986 & $-$11.934 \\
2000 & \textbf{74.224} & \textbf{74.332} & \textbf{11.709} & $-$11.658 \\
5000 & 74.181 & 74.255 & 11.667 & $-$11.616 \\
Dual ($\lambda_{\mathrm{init}}=500$) & 73.881 & 74.207 & 12.039 & $-$11.988 \\
\bottomrule
\end{tabular}
\end{table}

\section{Implementation Details}
\label{sec:impl}

\textbf{GAM optimizer.}
We use the GAM implementation with SGD base optimizer.
The perturbation radius is $\rho=0.05$ (gradient-norm perturbation),
with gradient norm perturbation radius $0.1$.
The GAM closure strictly evaluates $\Ltask$ only, with no Fisher penalty included.

\textbf{Fisher computation.}
Fisher is computed as the empirical Fisher over up to \texttt{ewc\_max\_batches}=100 mini-batches
at the end of each task.
For each mini-batch, we compute $\nabla_{\DeltaW_l} \log p(y|x;\phi)$ via backpropagation
and accumulate the squared norm.
The Fisher is then trace-normalized (Eq.~\eqref{eq:trace_norm}) with $\epsilon=10^{-8}$.

\textbf{Snapshot management.}
At the end of task $t$, we save the snapshot $\DeltaW^*_t = \{B_l A_l\}$ for all 24 adapter pairs.
The EWC penalty accumulates over all past tasks $t' < t$ with a decay factor $\gamma=0.9$
applied to the Fisher matrices.

\textbf{Learning rate schedule (Track~1, fine-grained, $r=16$).}
Cosine decay, $\mathrm{lr}=0.01$ uniform; 50 epochs for ImageNet-R, 20 epochs for all other datasets.

\textbf{Learning rate schedule (Track~2, PECL benchmarks).}
Epochs per task: CIFAR-100 20, ImageNet-A 20, ImageNet-R 50, DomainNet 5.
Rank: CIFAR-100 / ImageNet-A / ImageNet-R use $r=10$; DomainNet uses $r=30$.
The \texttt{config\_Expand} suite explores four learning rate strategies per dataset:
(a)~\emph{Uniform baseline}: $\mathrm{init\_lr}=\mathrm{lr}=0.02$;
(b)~\emph{Warmup-decay}: $\mathrm{init\_lr}\in\{0.02$--$0.04\}$ with $\mathrm{lr}=0.015$--$0.02$;
(c)~\emph{Equal-high}: $\mathrm{init\_lr}=\mathrm{lr}\in\{0.03$--$0.05\}$.
The best per-dataset variant is reported as the final Track~2 result.

\textbf{Gradient logging.}
For mechanistic analysis, we log $\cosff$, fisher/clean ratio, flat/clean ratio,
and \texttt{ewc\_penalty\_value} at the end of each task.
These statistics are written to the JSON output file for each run.
